\section{The binomial example and the shape of the likelihood}
\label{sec:binomial-shape}

Consider the binomial model with probability of success $q$ and $n$ known. The
probability mass function of the count $\Upsilon$ is
\begin{equation}
  \dens_\Upsilon(v; q, n) = \Prob(\Upsilon = v \given n, q)
  = \binom{n}{v} q^v (1-q)^{n-v},
  \qquad v \in \{0,1,\dots,n\}.
  \label{eq:pmf}
\end{equation}
Fixing an observed count $v$ and reading \eqref{eq:pmf} as a function of $q$ gives
the likelihood
\begin{equation}
  \like(q \given v, n) = \binom{n}{v} q^v (1-q)^{n-v},
  \qquad q \in [0,1].
  \label{eq:lik}
\end{equation}

\paragraph{Why it has that shape.}
As a function of $q$, expression \eqref{eq:lik} is the kernel of a
$\mathrm{Beta}(v+1,\, n-v+1)$ density; everything about its shape follows from
that. The binomial coefficient is constant in $q$: it rescales vertically and
plays no role in location or shape. The shape is the product of two competing
power terms --- $q^v$ rewards larger $q$ (pulling mass toward $1$), while
$(1-q)^{n-v}$ rewards smaller $q$ (pulling toward $0$) --- and they balance at
an interior point.

Work with the log-likelihood
\begin{equation}
  \loglike(q) = \text{const} + v\log q + (n-v)\log(1-q).
\end{equation}
Then
\begin{equation}
  \loglike'(q) = \frac{v}{q} - \frac{n-v}{1-q} = 0
  \;\Longrightarrow\; v(1-q) = (n-v)q
  \;\Longrightarrow\; q = \frac{v}{n},
\end{equation}
which is exactly the MLE, and
\begin{equation}
  \loglike''(q) = -\frac{v}{q^2} - \frac{n-v}{(1-q)^2} < 0
  \quad \text{on } (0,1),
\end{equation}
so $\loglike$ is strictly concave. Hence $\like$ is strictly log-concave,
therefore unimodal with a single peak at $q = v/n$.

For $0 < v < n$ the endpoints vanish, $\like(0) = \like(1) = 0$, giving the
characteristic single bump: rising from $0$, peaking at $v/n$, decaying back to
$0$. The degenerate cases collapse the bump against a wall: $v = 0$ gives
$(1-q)^n$, monotone decreasing with mode at $q = 0$; $v = n$ gives $q^n$,
monotone increasing with mode at $q = 1$.

\begin{remark}
$\like$ is smooth (a degree-$n$ polynomial in $q$), which is what is meant by
calling it ``continuous.'' It is \emph{not} a density in $q$: it does not
integrate to $1$ over $q$. Normalizing requires dividing by
$\Betafn(v+1, n-v+1)$, and the fact that this normalization produces the Beta
distribution is precisely Beta--Binomial conjugacy.
\end{remark}
